\chapter{Gradient Descent and Training}
Gradient descent (GD) is the most common method of training neural networks. GD makes a very simple approach. Given a loss function $L\bra{\boldsymbol{\Theta}}$ that depends on the model parameters $\boldsymbol{\Theta}$, we assume that a step in the negative direction of the gradient w.r.t. the parameters is a step towards the minimum.
\begin{align}
	\label{eq:gradient-descent}
	\boldsymbol{\Theta}_{k+1}&=\boldsymbol{\Theta}_k - \eta\cdot\nabla L
\end{align}
where $\eta$ is called the learning rate.\par
This method has some problems in general and some that have to do with neural networks. 
\textbf{Problems with standard GD}
\begin{itemize}
	\item Learning rate $\eta$ is a hyperparameter that might not fit the problem well
	\item Works only if in the vicinity of the global minimum
	\item Even if in a convex region around the global minimum, might overshoot
\end{itemize}
For those reasons, standard gradient descent is almost unfeasible for training. Instead of tweaking this formular with ever compounding non-justifiable modifications, we should start out with what we know about the problem and what method works best in general.\par
The desired method is a Trust-Region Reflective Algorithm like Levenberg-Marquardt and we need to take into account that we have no method of good initialization and that the loss function will be riddled with local extrema.
\section{Newton-Raphson Optimization}
Let's assume a function $L:\mathbb{R}^n\to\mathbb{R}$ with $\boldsymbol{x}\mapsto L$. Then we can write the second order approximation around a point $\boldsymbol{x}_0$ with a truncated Taylor series.
\begin{align}
	\label{eq:taylor2}
	L_2\bra{\boldsymbol{x}}&= L\bra{\boldsymbol{x}_0}+\nabla L\bra{\boldsymbol{x}_0}\cdot\bra{\boldsymbol{x-x}_0}+\frac{1}{2}\bra{\boldsymbol{x-x}_0}\boldsymbol{H}\bra{\boldsymbol{x}_0}\bra{\boldsymbol{x-x}_0}
\end{align}
Note: The gradient is a row vector. The usage of the partial derivatives as a column vector is only possible in orthonormal spaces. The correct gradient transform behaviour requires it to be a row when using simplified notation instead of component equations.\par
What the NR-Method does is solving a proxy problem. We have no method to solve a general nonlinear system of equations or finding the minimum. Therefore, we approximate the problem with a paraboloid, a second order polynomial expansion and solve for its minimum, which is trivial, because we just need to solve the first derivatives equals zero.\par
\begin{align}
	\label{eq:taylor22}
	\nabla L_2&=\nabla L\bra{\boldsymbol{x}_0}+\boldsymbol{H}\bra{\boldsymbol{x}_0}\bra{\boldsymbol{x-x}_0}\overset{!}{=}\boldsymbol{0}\\
	\boldsymbol{x}&=\boldsymbol{x}_0-\boldsymbol{H}^{-1}\bra{\boldsymbol{x}_0}\nabla L \bra{\boldsymbol{x}_0}
\end{align}
If we now make this an iteration
\begin{align}
	\label{eq:nr}
	\boldsymbol{x}_{n+1}&=\boldsymbol{x}_n-\boldsymbol{H}^{-1}\bra{\boldsymbol{x}_n}\nabla L \bra{\boldsymbol{x}_n}
\end{align}
we obtain a formula that looks like gradient descent but with some sort of \enquote{scaling} of the gradient that depends on the current iteration step.
\subsubsection{Discussion of the method}
The Newton-Raphson method is mathematically elegant but fragile in practice. Its primary vulnerability lies in the \textbf{assumption of local quadraticity}. If the second-order Taylor expansion does not sufficiently reflect the true curvature of the loss surface, the iteration step becomes unreliable. 

Key structural weaknesses include:
\begin{itemize}
	\item \textbf{Topographical Blindness}: The optimizer assumes a \enquote{well-behaved} landscape. In non-convex regions, the Hessian $\boldsymbol{H}$ may not be positive definite. In such cases, the method might move toward a local maximum or a saddle point rather than a minimum.
	\item \textbf{The Convergence Radius}: Convergence is only guaranteed if the initial point $\boldsymbol{x}_0$ lies within a strict convex basin of the minimum. Outside this region, the lack of a global \enquote{map} often leads to erratic updates.
	\item \textbf{Overshooting and Divergence}: Even within a convex region, if the curvature is nearly flat (small eigenvalues in the Hessian), the term $\boldsymbol{H}^{-1}$ produces an aggressively large step that can catapult the iteration out of the stable region entirely.
\end{itemize}

To mitigate this, Levenberg and Marquardt introduced the concept of \textbf{Trust Regions}. Rather than taking the full NR-step blindly, the method constrains the update to a neighborhood where the quadratic approximation is \enquote{trusted}. Mathematically, this is implemented by augmenting the Hessian: $(\boldsymbol{H} + \lambda \boldsymbol{I})$. 

This damping factor $\lambda$ performs a dual role:
\begin{enumerate}
	\item \textbf{Interpolation}: It smoothly interpolates between a cautious Gradient Descent step (large $\lambda$) and a precise Newton step (small $\lambda$).
	\item \textbf{Regularization}: It ensures the matrix is invertible and positive definite, effectively \enquote{warping} the local topography to be more spherical and stable.
\end{enumerate}

While the explicit Trust Region logic of varying $\lambda$ based on gain-ratio is rarely used in deep learning due to the overhead of Hessian computation, the underlying principle of \textbf{adaptive damping} is the bridge toward more modern, scalable optimizers.
\section{Handling the Hessian}
Since the Hessian is unfeasible to construct with a high number of parameters but multiplying with the inverse hessian ensures that the step is taken according to a quadratic approximation rather than just a linear one, we need a way to approximate the Hessian. For this, we introduce the Likelihood function. 
\subsection{Statistics intermission}
To derive a feasible approximation of the Hessian, we introduce the joint probability density function (PDF) of a sample. Consider $N$ individual observations where the PDF of the $k$-th sample is expressed as $p\bra{x_k, \boldsymbol{\Theta}}$. Under the assumption of identically and independently distributed (iid) observations, the joint PDF is defined as:
\begin{align}
	\label{eq:pdf}
	P\bra{\boldsymbol{x}, \boldsymbol{\Theta}}&=\prod\limits_{k=1}^{N}p\bra{x_k, \boldsymbol{\Theta}}
\end{align}
Here, the bold-faced $\boldsymbol{x}$ represents the set of all observations, and $\boldsymbol{\Theta}$ denotes the distribution parameters. Conventionally, the joint PDF is interpreted as the probability density of a sample $\boldsymbol{x}$ given a fixed set of parameters $\boldsymbol{\Theta}$.\par
However, in an optimization context, this perspective is inverted. We are not interested in the probability of a sample we already possess; rather, we seek the parameters $\boldsymbol{\Theta}$ that are \textbf{most likely} to have generated the observed data. This shift in viewpoint is formalized by the likelihood function:
\begin{align}
	\label{eq:likelihood}
	\mathcal{L}\bra{\boldsymbol{\Theta} \left| \boldsymbol{x} \right.}&=\prod\limits_{k=1}^{N}p\bra{x_k, \boldsymbol{\Theta}}
\end{align}
While the functional form remains unchanged, $\mathcal{L}$ is now interpreted as the likelihood of the parameters $\boldsymbol{\Theta}$ given the evidence $\boldsymbol{x}$.

\subsection{Connection to the loss function}
To proceed, we must assume a specific functional form for the distribution. We assume a normal distribution—not merely for its ubiquity, but because the majority of frequentist analyses become mathematically intractable without the guarantee of normally distributed residuals. From a pragmatic standpoint, as long as the data is roughly symmetric and unimodal, the Gaussian assumption yields usable results. Under this assumption, the marginal PDF is:
\begin{align}
	\label{eq:norma}
	p\bra{x_k, \mu, \sigma}&=\frac{1}{\sqrt{2\pi}\sigma}\cdot\exp\bra{-\frac{1}{2}\bra{\frac{x_k-\mu}{\sigma}}^2}
\end{align}
Substituting this into the likelihood function for the entire sample yields:
\begin{align}
	\mathcal{L}(\mu,\sigma \left|\, \boldsymbol{x}\right.)&=\frac{1}{\bra{\sqrt{2\pi}\sigma}^N}\cdot\exp\bra{-\frac{1}{2\sigma^2}\sum\limits_{k=1}^N\bra{x_k-\mu}^2}
\end{align}
To simplify the optimization, we apply the logarithm to the likelihood function. Since the logarithm is a monotonic transformation, the location of the extremum remains unchanged:
\begin{align}
	\label{eq:loglikelihood}
	\log\mathcal{L}&=-N \log\bra{\sqrt{2\pi}\sigma}-\frac{1}{2\sigma^2}\underset{L}{\underbrace{\sum\limits_{k=1}^N\bra{x_k-\mu}^2}}
\end{align}
This reveals the fundamental link between statistics and neural network training: the negative log-likelihood of a Gaussian distribution is equivalent to the quadratic loss. By treating the network output as the parameter $\mu$ and observing that the scaling factor $1/\sigma^2$ and the constant offset do not shift the minimum, we can establish the following relationship without loss of generality:
\begin{align}
	\label{eq:losslikelihood}
	L \propto -\log\mathcal{L}
\end{align}

\subsection{Analyzing the Hessian}
We now examine the structure of the Hessian to find a suitable approximation for the Newton step. Given the relationship $L = c \log \mathcal{L}$, the Hessian is:
\begin{align}
	\label{eq:hessian-derivation}
	\boldsymbol{H} &= -c \frac{\partial}{\partial \boldsymbol{\Theta}} \left( \frac{\partial \log \mathcal{L}}{\partial \boldsymbol{\Theta}} \right) \\
	&= -c \frac{\partial}{\partial \boldsymbol{\Theta}} \left( \frac{1}{\mathcal{L}} \nabla \mathcal{L} \right) \\
	&= c \left[ \frac{1}{\mathcal{L}^2} \nabla \mathcal{L}^\top \nabla \mathcal{L} - \frac{1}{\mathcal{L}} \nabla^2 \mathcal{L} \right]
\end{align}
By substituting the identity $\frac{\nabla \mathcal{L}}{\mathcal{L}} = \nabla \log \mathcal{L}$, we obtain:
\begin{align}
	\label{eq:hessian-expanded}
	\boldsymbol{H} &= c \bra{\nabla \log \mathcal{L}}^\top \bra{\nabla \log \mathcal{L}} - \frac{c}{\mathcal{L}} \nabla^2 \mathcal{L}
\end{align}
To arrive at a generalizable optimizer, we must consider the \textbf{expected Hessian} over the data distribution. The expectation of the second term in equation \eqref{eq:hessian-expanded} can be evaluated as follows:
\begin{align}
	E\left[ \frac{c}{\mathcal{L}} \nabla^2 \mathcal{L} \right] &= \int \frac{c}{\mathcal{L}(\boldsymbol{x})} \nabla^2 \mathcal{L}(\boldsymbol{x}) \cdot \mathcal{L}(\boldsymbol{x}) \, d\boldsymbol{x} \\
	&= c \int \nabla^2 \mathcal{L}(\boldsymbol{x}) \, d\boldsymbol{x}
\end{align}
Under the assumption that the likelihood function is sufficiently smooth to allow the swap of the integral and the Laplace operator (Leibniz's rule), we invoke the conservation of total probability:
\begin{align}
	c \nabla^2 \int \mathcal{L}(\boldsymbol{x}) \, d\boldsymbol{x} = c \nabla^2 (1) = 0
\end{align}
Consequently, the second term vanishes in expectation. This yields the fundamental identity:
\begin{align}
	\label{eq:fisher-identity}
	E[\boldsymbol{H}] &= c \cdot E\left[ \bra{\nabla \log \mathcal{L}}^\top \bra{\nabla \log \mathcal{L}} \right]
\end{align}
By recalling $\nabla L = -c \nabla \log \mathcal{L}$, we can express the expected Hessian solely through the first-order gradients of our loss:
\begin{align}
	E[\boldsymbol{H}] &= \frac{1}{c} E\left[ \nabla L^\top \nabla L \right]
\end{align}
This result is profound: it demonstrates that while the true Hessian is computationally inaccessible, we can rigorously approximate the expected curvature of the loss surface using only the outer product of the gradients, provided we average over a sufficient number of observations. We recall that the proportionality constant $c=1/\sigma^2$. This is important later, because in ADAMprop the scaling factor of the learning rate is not accidently called the variance.
\section{The Mechanical Necessity of Expected Values}
If we take a look at classical regression problems, we usually define a quadratic loss over a dataset $\mathcal{D}$ as follows:
\begin{align}
	L = \sum\limits_k \left( y_k - \hat{y}_k(\boldsymbol{\Theta}) \right)^2
\end{align}
A step of ordinary gradient descent is then:
\begin{align}
	\boldsymbol{\Theta}_{n+1} = \boldsymbol{\Theta}_{n} - \eta \cdot \nabla L
\end{align}
However, this is often impractical for training neural networks. Computing the full gradient requires the entire dataset to be in memory, and updating parameters only after a full pass is computationally inefficient. To address this, we recognize that the gradient operator is linear and can be moved inside the summation:
\begin{align}
	\nabla L = \sum\limits_k \nabla L_k(\boldsymbol{\Theta}_n)
\end{align}
While we can calculate contributions $\nabla L_k$ observation-wise, waiting until the end of the dataset to move is \enquote{thinking a lot and moving once.} To \enquote{move while thinking,} we can apply gradient descent to a single observation:
\begin{align}
	\boldsymbol{\Theta}_{k+1} = \boldsymbol{\Theta}_{k} - \frac{\eta}{N} \nabla L_k(\boldsymbol{\Theta}_k)
\end{align}
Note that the learning rate is scaled by $N$ because we are now taking $N$ steps per epoch. If we iterate this process for the full dataset, the final state becomes:
\begin{align}
	\boldsymbol{\Theta}_{N} = \boldsymbol{\Theta}_{0} - \sum\limits_{k=1}^N \frac{\eta}{N} \nabla L_k(\boldsymbol{\Theta}_k) = \boldsymbol{\Theta}_{0} - \eta \cdot \frac{1}{N} \sum\limits_{k=1}^N \nabla L_k(\boldsymbol{\Theta}_k)
\end{align}
This expression identifies the full update as a function of the \textbf{arithmetic mean} of the local gradient contributions. Since the arithmetic mean is an unbiased estimator for the expected value, we can reformulate the total update as:
\begin{align}
	\boldsymbol{\Theta}_{N} = \boldsymbol{\Theta}_0 - \eta \cdot \mathbb{E}\left[ \nabla L_i \right]
\end{align}
Here, $i$ represents a random index. This motivates an update law where we replace the noisy, sampled gradient with its expected value:
\begin{align}
	\boldsymbol{\Theta}_{k+1} = \boldsymbol{\Theta}_k - \eta \cdot \mathbb{E}\left[ \nabla L_k \right]
\end{align}

\section{The Noisy Integrator}
The step-wise gradient contributions act as a noisy input to an integrator. Consider the simplified system:
\begin{align}
	y_{k+1} = y_k + \epsilon_k
\end{align}
If $\epsilon_k$ is a random variable with expectation $\mu$, the output $y_N$ after $N$ steps approximates $N \cdot \mu$. A mean value answers the question: \enquote{With what single value could I replace every observation to achieve the same result?} For an additive system, this is the arithmetic mean. By replacing the noisy contribution $\epsilon_k$ with the expected value $\mu$, we obtain a smooth integrator:
\begin{align}
	y_{k+1} = y_k + \mu
\end{align}
Since $\mu$ is unknown and computing the full mean every step is expensive, we employ a recursive estimator—a moving average or low-pass filter:
\begin{align}
	\label{eq:update}
	m_{k+1} = \alpha \cdot m_k + (1-\alpha) \cdot \epsilon_k
\end{align}
With $\alpha \approx 1$, we prioritize the accumulated history over the current noisy measurement, effectively suppressing high-frequency noise. In the context of learning, where the \enquote{noise} $\nabla L_k$ also depends on the current state $\boldsymbol{\Theta}_k$, such filtering is critical to prevent localized gradients from derailing the global optimization path.
\section{Respecting Cold Starts}
There is one part left that is usually not treated in standard filtering: the transient response. In most system theory approaches, we focus on the steady-state response, which is somewhat equivalent to looking at the expectation of the series. First, we ground the iteration in an initial state to resolve the recurrence:
\begin{align}
	m_2 &= \alpha\cdot m_1 + \bra{1-\alpha}\cdot\epsilon_1 \\
	m_3 &= \alpha\cdot m_2 + \bra{1-\alpha}\cdot\epsilon_2 \\
	&= \alpha^2\cdot m_1 + \bra{1-\alpha}\cdot\bra{\alpha\cdot\epsilon_1+\epsilon_2}
\end{align}
The pattern emerges immediately, and the closed-form expression is given by:
\begin{align}
	\label{eq:noisy-integrator-closed}
	m_{k+1} = \alpha^k m_1 + \bra{1-\alpha} \sum_{i=1}^k \epsilon_i \cdot \alpha^{k-i}
\end{align}
In the steady state, the contribution of the initial state $m_1$ approaches zero. Taking the expectation of \eqref{eq:noisy-integrator-closed} yields:
\begin{align}
	\label{eq:expected-noise}
	\mathbb{E}\bra{\bra{1-\alpha} \sum_{i=1}^k \epsilon_i \cdot \alpha^{k-i}} 
	&= \bra{1-\alpha} \sum_{i=1}^k \mathbb{E}\bra{\epsilon_i} \cdot \alpha^{k-i} \\
	&= \mu \cdot \bra{1-\alpha} \sum_{i=1}^k \alpha^{k-i}
\end{align}

\subsubsection{Resolution of the sum}
\begin{enumerate}
	\item Substitute $j=k-i$ in $\sum_{i=1}^k \alpha^{k-i} \Rightarrow \sum_{j=k-1}^0 \alpha^j = \sum_{j=0}^{k-1} \alpha^j = S_{k-1}$.
	\item Multiply $S_{k-1}$ by $\alpha$ and subtract it from the original sum to isolate the geometric boundary:
	\begin{align}
		\bra{1-\alpha} S_{k-1} &= \sum_{j=0}^{k-1} \alpha^j - \sum_{j=0}^{k-1} \alpha^{j+1} \\
		&= \sum_{j=0}^{k-1} \alpha^j - \sum_{j=1}^{k} \alpha^{j} \\
		&= \bra{1 + \sum_{j=1}^{k-1} \alpha^j} - \bra{\alpha^k + \sum_{j=1}^{k-1} \alpha^j} \\
		&= 1 - \alpha^k
	\end{align}
	\item Thus, $S_{k-1} = \frac{1-\alpha^k}{1-\alpha}$.
\end{enumerate}

This leads to the biased expectation:
\begin{align}
	\mathbb{E}\bra{\bra{1-\alpha} \sum_{i=1}^k \epsilon_i \cdot \alpha^{k-i}} &= \mu \cdot \bra{1 - \alpha^k}
\end{align}
For rising $k$, the expected value of the low-pass filter approaches the true expectation of the noise. However, if the system is initialized with $m_1=0$, this convergence in the transient phase is sluggish. To optimize for \enquote{cold starts}, we tweak the update by dividing by the bias-correction term $\bra{1-\alpha^k}$:
\begin{align}
	\hat{m}_{k+1} &= \frac{m_{k+1}}{1-\alpha^k} = \frac{\alpha}{1-\alpha^k} m_k + \frac{1-\alpha}{1-\alpha^k} \epsilon_k
\end{align}
At the cold start ($k=1$), we find $\hat{m}_2 = \epsilon_1$, effectively ignoring the zero-initialized memory. As $k$ grows, this transient correction factor approaches unity. Crucially, this correction is a \enquote{multiplication after the fact}---the correction is applied only during the descent step and is not fed back into the recursive state of the filter.
\section{Mini-Batches: Consistency of the Estimator}
The derivation for the expected value in equation \eqref{eq:expected-noise} was demonstrated for single-observation updates ($B=1$). In practice, we aggregate $B$ observations into a \enquote{mini-batch} to utilize the parallel processing capabilities of modern hardware. This shift is not a change in the underlying optimization theory, but a change in the variance of our estimator. By the linearity of the expectation operator, the expected value of the batch mean remains identical to the population mean $\mu$:
\begin{align}
	\mathbb{E}\bra{\nabla L_{batch}} = \mathbb{E}\bra{\frac{1}{B} \sum_{i=1}^B \nabla L_i} = \frac{1}{B} \sum_{i=1}^B \mathbb{E}\bra{\nabla L_i} = \mu
\end{align}
The mini-batch simply provides a lower-variance sample of the gradient at each step. Using a mini-batch allows us to maintain a high \enquote{refresh rate} for the parameters $\boldsymbol{\Theta}$ while benefiting from the central limit theorem—the noise floor of the gradient is suppressed by a factor of $\sqrt{B}$ before it even reaches our recursive filter.

\section{Numerical Tweaks and Dimensional Sanity}
In the transition from the theoretical Newton-Raphson step to a robust algorithm, two adjustments are made that often appear heuristic but follow a distinct internal logic. 

First, we address the division by the second moment. While equation \eqref{eq:fisher-identity} links the Hessian to the variance, empirical evidence shows that scaling the learning rate by the \textbf{standard deviation} (the square root of the second moment) yields more stable convergence. This can be viewed as a dimensional necessity: if we consider the units of the parameter $\boldsymbol{\Theta}$ to be $[U]$, the gradient $\nabla L$ has units $[U^{-1}]$. A true Newton step involves $\boldsymbol{H}^{-1}\nabla L$, which has units $[U^2 \cdot U^{-1}] = [U]$. By using the square root of the second moment (the RMS of the gradients), we maintain this dimensional consistency.

Second, we introduce the smoothing term $\epsilon$:
\begin{align}
	\Delta \boldsymbol{\Theta} = - \eta \cdot \frac{\hat{m}}{\sqrt{\hat{v}} + \epsilon}
\end{align}
This is not merely a \enquote{hack} to avoid division by zero. From a control perspective, $\epsilon$ represents a \textbf{damping floor}. It prevents the gain of the optimizer from approaching infinity in regions where the loss surface is perfectly flat (where the gradient and its variance vanish), effectively enforcing a maximum allowable step size.

\section{Construction of the Adamprop Algorithm}
We can now synthesize the full algorithm.

\subsection*{Starting Point}
We start out with the Newton-Raphson update formula for the full dataset loss $L$:
\begin{align}
	\boldsymbol{\Theta}_{n+1}&=\boldsymbol{\Theta}_{n} - \boldsymbol{H}^{-1}\bra{\boldsymbol{\Theta}_n}\cdot\nabla L\bra{\boldsymbol{\Theta}_n}
\end{align}

\subsection*{Observation-wise Update}
We make the update every observation. To avoid stochastic effects, we replace the current values of the Hessian and gradient with the expected values.
\begin{align}
	\boldsymbol{\Theta}_{k+1}&=\boldsymbol{\Theta}_{k} - \mathbb{E}\bra{\boldsymbol{H}_k\bra{\boldsymbol{\Theta}_k}}^{-1}\cdot\mathbb{E}\bra{\nabla L_k\bra{\boldsymbol{\Theta}_k}}
\end{align}
For formula sanitization, we drop the arguments and indices of the Hessian and gradient:
\begin{align}
	\boldsymbol{\Theta}_{k+1}&=\boldsymbol{\Theta}_{k} - \mathbb{E}\bra{\boldsymbol{H}}^{-1}\cdot\mathbb{E}\bra{\nabla L}
\end{align}

\subsection*{Using the FIM-Equality}
We use the Fisher Information Matrix (FIM) identity to get rid of the Hessian. Note that since the gradient is a row vector, the dyadic product is formed by $\nabla L^\top \nabla L$.
\begin{align}
	\boldsymbol{\Theta}_{k+1}&=\boldsymbol{\Theta}_{k} - \frac{1}{\sigma^2}\mathbb{E}\bra{\boldsymbol{\nabla L^\top\nabla L}}^{-1}\cdot\mathbb{E}\bra{\nabla L}
\end{align}

\subsection*{Approximation of the dyadic gradient product}
Due to computational demand, we assume that the main diagonal of the dyadic gradient product is a feasible approximation.
\begin{align}
	\boldsymbol{\Theta}_{k+1}&=\boldsymbol{\Theta}_{k} - \frac{1}{\sigma^2}\mathbb{E}\bra{\mathrm{diag}\bra{g_i^2}}^{-1}\cdot\mathbb{E}\bra{\nabla L}
\end{align}

\subsection*{Running estimation of the expected values}
We introduce two low-pass filters for the gradient and the gradient squared. We do it component-wise for the diagonal matrix.
\begin{align}
	\boldsymbol{m}_{k+1}&=\alpha\cdot \boldsymbol{m}_{k} + \bra{1-\alpha}\cdot \nabla L_k\\
	\boldsymbol{v}_{k+1}&=\beta\cdot \boldsymbol{v}_{k} + \bra{1-\beta}\cdot \nabla L_k\odot\nabla L_k
\end{align}
$\alpha$ and $\beta$ are constants approaching $1$ from below. The closer to one, the lower the cut-off frequency and the slower the estimation. Usually $\beta$ is chosen even more conservatively than $\alpha$. The nomenclature with $\boldsymbol{m}$ and $\boldsymbol{v}$ comes from the terms \enquote{momentum} and \enquote{variance}. The Fisher Information Matrix is an approximation of the covariance matrix of the estimator, therefore the name variance fits. The running estimate of the gradient is called momentum, because the running estimate cannot be changed erratically and a directional change is not instantaneous. We use the Hadamard product for the gradient squared here.
\begin{align}
	\boldsymbol{\Theta}_{k+1}&=\boldsymbol{\Theta}_{k} - \frac{1}{\sigma^2}\frac{1}{\boldsymbol{v}_k}\odot\boldsymbol{m}_k
\end{align}
The division by the vector indicates an element-wise reciprocal.

\subsection*{Consider Cold Starts}
Replace the variance and momentum by their corrections to account for initialization bias:
\begin{align}
	\hat{\boldsymbol{m}}_k&=\frac{\boldsymbol{m}_k}{1-\alpha^k}\\
	\hat{\boldsymbol{v}}_k&=\frac{\boldsymbol{v}_k}{1-\beta^k}
\end{align}

\subsection*{Numerical Tweaks for Sanitization}

\subsection*{Numerical Tweaks for Sanitization}

\subsubsection*{Learning Rate as a Safety Valve}
The jump from a clean Newton-Raphson update to a messy stochastic one is a recipe for instability. In a perfect world, NR tells you exactly how far to go, but in our world, without trust regions or line-search, the algorithm is prone to overshooting and "exploding" even when the path looks clear. Plus, we don't actually know the noise variance $\sigma^2$ from the FIM. So, we simplify: we drop the theoretical $1/\sigma^2$ and swap in an empirical learning rate $\eta$. This acts as our global damping factor, letting us manually dial back the update magnitude so the whole thing doesn't vibrate itself to pieces during convergence.

\subsubsection*{Standard Deviation over Variance}
The raw NR-derivation tells us to divide by the variance (the second moment). In practice, that is way too aggressive for a stochastic gradient field; it scales the steps too wildly when the noise kicks in. It turns out that taking the square root—normalizing by the \textit{standard deviation} instead—is much more robust. This basically turns our update into a Signal-to-Noise Ratio (SNR) weight. Components with high variance (lots of noise) get damped down, while components with a consistent, low-variance signal get the green light. It’s an automated way of preconditioning the manifold without the instability of a full second-order scaling.

\subsubsection*{Preventing Singularities}
When you hit a flat region in the loss landscape, your gradient power $\boldsymbol{v}_k$ effectively vanishes. Mathematically, you're looking at a division-by-zero disaster. To keep the code from crashing and the updates from hitting infinity, we throw in a tiny regularization constant $\epsilon$. This "floor" ensures the denominator never hits zero, keeping the update well-conditioned even in the flattest parts of the space. It also serves as a cap on our maximum step size, so we don't go flying off into space just because the local gradient is near zero.

\subsection*{Adamprop Update Formula}
Integrating the recursive estimators with their respective bias corrections, we arrive at the final update law:
\begin{align}
	\boldsymbol{\Theta}_{k+1} &= \boldsymbol{\Theta}_{k} - \frac{\eta}{\sqrt{\hat{\boldsymbol{v}}_k} + \epsilon} \odot \hat{\boldsymbol{m}}_k
\end{align}
\section{The Training Algorithm}
\begin{algorithm}[H]
	\SetAlgoLined
	\DontPrintSemicolon
	\caption{Full Training Algorithm on mini-batches}
	
	\KwIn{Training set $\mathcal{D}$, Batch size $N$, Learning rate $\eta$, Scaling factors $\alpha$ and $\beta$}
	\BlankLine
	\textbf{Initialize} momentum $\boldsymbol{m}_1=0$\;
	\textbf{Initialize} variance $\boldsymbol{v}_1=0$\;
	\textbf{Set} number of epochs $N_\mathrm{epoch}$\;
	\BlankLine
	\While{$\mathrm{epoch}<N_\mathrm{epoch}$}{
		\textbf{Shuffle} dataset $\mathcal{D}$\;
		\ForEach{batch $\mathcal{B} \subset \mathcal{D}$ of size $N$}{
			\tcc{Parallelization over samples $k \in \mathcal{B}$}
			\ForEach{sample $k$}{
				1. \textbf{Forward Pass}: Compute and store activations $a^l_i$\;
				2. \textbf{Backward Pass}: Compute and store error signals $\delta^l_i$\;
				3. \textbf{Local Gradient}: Calculate weight update contribution $\Delta w_k$\;
			}
			4. \textbf{Aggregation}: Sum contributions $\Delta W = \sum_{k=1}^{N} \Delta w_k$\;
			5. \textbf{Update}: Adjust weights $w \leftarrow w - \eta \Delta W$\;
		}
		\textbf{Check Termination}: Evaluate loss/convergence to decide if another epoch is needed\;
	}
\end{algorithm}
